\documentclass[11pt]{article}
% ======================================================================
%  Preambule commun - Sujets Bac Serie A (Mathematiques)
%  Style sobre : noir et blanc, sans couleurs, sans filigrane,
%  sans en-tetes ni pieds de page decoratifs.
% ======================================================================
\usepackage[utf8]{inputenc}
\usepackage[T1]{fontenc}
\usepackage{lmodern}
\usepackage{textcomp}
\usepackage{amsmath,amssymb}
\usepackage{enumitem}
\usepackage{array}
\usepackage[a4paper,margin=2.3cm]{geometry}
\usepackage{fancyhdr}
\usepackage{tikz}

% Symbole degre robuste + justification souple
\DeclareUnicodeCharacter{00B0}{\ensuremath{^\circ}}
\tolerance=2000
\emergencystretch=2em
\frenchspacing

% En-tete / pied de page minimalistes (numero de page seul)
\pagestyle{fancy}
\fancyhf{}
\renewcommand{\headrulewidth}{0pt}
\fancyfoot[C]{\thepage}

% Macros mathematiques
\newcommand{\R}{\mathbb{R}}
\newcommand{\N}{\mathbb{N}}
\newcommand{\Z}{\mathbb{Z}}
\newcommand{\vi}{\vec{\imath}}
\newcommand{\vj}{\vec{\jmath}}

% Titres de blocs
\newcommand{\exo}[2]{\bigskip\noindent{\large\bfseries Exercice #1}\hfill\textit{#2}\par\nopagebreak\smallskip}
\newcommand{\partie}[1]{\medskip\noindent{\bfseries Partie #1}\par\nopagebreak\smallskip}
\newcommand{\entetesujet}[1]{\begin{center}{\Large\bfseries #1}\end{center}\vspace{0.3em}\hrule\vspace{1.1em}}

% Questions numerotees : niveau 1 = chiffres gras, niveau 2 = a. b. c.
\newlist{questions}{enumerate}{1}
\setlist[questions]{label=\textbf{\arabic*.},leftmargin=1.8em,itemsep=3pt,topsep=3pt}
\newlist{sousq}{enumerate}{1}
\setlist[sousq]{label=\textbf{\alph*.},leftmargin=1.8em,itemsep=2pt,topsep=2pt}


\begin{document}

\entetesujet{Sujet bac 2016 -- Série A}

\exo{1}{8 points}
On donne les nombres $A = 2\cdot 10^{-4} + 103\cdot 10^{-2}$ et $B = 301\cdot 10^{-3} \times 11\cdot 10^{-2}$.
\begin{questions}
  \item Effectuer les calculs des nombres $A$ et $B$. On donnera les réponses sans puissances de 10.
  \item On donne : $P(x) = -x^3 + 2x^2 + x - 2$.
  \begin{sousq}
    \item Calculer $P(1)$.
    \item Déterminer les réels $a$ et $b$ tels que $P(x) = (x - 1)(-x^2 + ax + b)$.
    \item Résoudre l'équation : $-x^2 + x + 2 = 0$. En déduire les solutions de $P(x) = 0$.
    \item En déduire les solutions de l'équation $-e^{3x} + 2\,e^{2x} + e^{x} - 2 = 0$.
  \end{sousq}
  \item Résoudre dans $\R^2$ le système d'équations suivant : $S_1 \begin{cases} x - y = 2 \\ 2x + 3y = -1 \end{cases}$

  En déduire les solutions du système suivant : $S_2 \begin{cases} \ln x - \ln y = 2 \\ 2\ln x + 3\ln y = -1 \end{cases}$
\end{questions}

\exo{2}{8 points}
On considère la fonction numérique $f$ à variable réelle $x$ définie par :
\[ f(x) = \frac{2 - x}{x - 1} \]
On désigne par $(\mathcal{C})$ la courbe représentative de $f$ dans le repère orthonormé $(O,\vi,\vj)$ du plan. Unité graphique : 2 cm.
\begin{questions}
  \item Déterminer l'ensemble de définition de $f$.
  \item Déterminer les réels $a$ et $b$ tels que pour tout $x \ne 1$, on a : $f(x) = a + \dfrac{b}{x - 1}$.
  \item On donne $\lim\limits_{x\to-\infty} f(x) = -1$ et $\lim\limits_{x\to+\infty} f(x) = -1$. Calculer la limite de $f$ à gauche et à droite en $x = 1$.
  \item \begin{sousq}
    \item Calculer la dérivée $f'$ de $f$.
    \item Dresser le tableau de variation de $f$.
    \item Que représentent les droites d'équations $x = 1$ et $y = -1$ pour la courbe $(\mathcal{C})$ de $f$ ?
  \end{sousq}
  \item \begin{sousq}
    \item Compléter le tableau de valeurs ci-dessous :
    \begin{center}\renewcommand{\arraystretch}{1.5}
    \begin{tabular}{|c|c|c|c|c|c|c|c|}
    \hline
    $x$ & $-3$ & $-2$ & $-1$ & $0$ & $1$ & $2$ & $3$ \\
    \hline
    $y$ & & & & & & & \\
    \hline
    \end{tabular}
    \end{center}
    \item Construire, dans le repère $(O,\vi,\vj)$ du plan, la courbe $(\mathcal{C})$ de $f$ et les droites d'équations $x = 1$ ; $y = -1$.
  \end{sousq}
\end{questions}

\exo{3}{4 points}
Le tableau suivant donne le poids $x$ en grammes et la taille $y$ en centimètres d'une population donnée.
\begin{center}\renewcommand{\arraystretch}{1.4}
\begin{tabular}{|c|c|c|c|c|c|c|c|c|c|c|}
\hline
Poids $x$ & 10 & 25 & 40 & 50 & 55 & 60 & 65 & 70 & 75 & 80 \\
\hline
Taille $y$ & 11 & 20 & 35 & 45 & 50 & 53 & 60 & 63 & 73 & 75 \\
\hline
\end{tabular}
\end{center}
\begin{questions}
  \item Représenter le nuage des points dans le plan muni d'un repère orthogonal. 1 cm pour 10 g et 1 cm pour 10 cm.
  \item La série ci-dessus est divisée en deux sous-séries telles que :

  \underline{Sous-série A}
  \begin{center}\renewcommand{\arraystretch}{1.4}
  \begin{tabular}{|c|c|c|c|c|c|}
  \hline Poids $x$ & 10 & 25 & 40 & 50 & 55 \\\hline Taille $y$ & 11 & 20 & 35 & 45 & 50 \\\hline
  \end{tabular}\end{center}

  \underline{Sous-série B}
  \begin{center}\renewcommand{\arraystretch}{1.4}
  \begin{tabular}{|c|c|c|c|c|c|}
  \hline Poids $x$ & 60 & 65 & 70 & 75 & 80 \\\hline Taille $y$ & 53 & 60 & 63 & 73 & 75 \\\hline
  \end{tabular}\end{center}
  \begin{sousq}
    \item Calculer les coordonnées de $G_1$ et $G_2$ points moyens respectifs des sous-séries A et B.
    \item Placer les points $G_1$ et $G_2$, puis tracer la droite $(G_1G_2)$ dans le repère orthogonal précédent. Que représente cette droite pour la série ?
  \end{sousq}
  \item On suppose $G_1(36\,;32{,}2)$ et $G_2(70\,;64{,}8)$.
  \begin{sousq}
    \item Déterminer une équation cartésienne de la droite $(G_1G_2)$.
    \item À l'aide de l'équation de la droite $(G_1G_2)$ obtenue, estimer la taille d'une personne ayant un poids de 97 grammes.
  \end{sousq}
\end{questions}

\end{document}
